\chapter{Background}
\label{chap:background}

This chapter collects the general background needed across the thesis.
Its purpose is to introduce broad concepts and research directions that will be revisited in the later technical chapters, without turning into a notation chapter or duplicating paper-specific preliminaries.

\TODO{Introduce the high-level thesis context around trustworthy generative modeling, including the role of constraints such as privacy, safety, and controlled generation.}

\TODO{Introduce the basic privacy motivation and the role of differential privacy in releasing data-derived artifacts.}

\TODO{Summarize the goal of synthetic tabular data generation and why it matters in sensitive application domains.}

\TODO{Briefly position the main families of private tabular synthesis methods that form the context for Chapter~\ref{chap:privpgd}.}

\TODO{Introduce the optimal-transport viewpoint at a high level and explain why geometry-aware discrepancies are useful in generative modeling.}

\TODO{Give the high-level intuition for sliced Wasserstein distances and why they are computationally attractive in high-dimensional generative problems.}

\TODO{Add a short background discussion on particle-based optimization methods and why they provide a natural fit for direct dataset construction.}

\TODO{Introduce memorization and regurgitation as a practical issue in text and code generation.}

\TODO{Provide a broad overview of training-time, post-hoc, and inference-time approaches to reducing copyright risk.}

\TODO{Motivate inference-time fusion methods as the immediate background for Chapter~\ref{chap:cpfuse}.}

\TODO{Introduce the causal-inference motivation for combining randomized trials and observational studies.}

\TODO{Explain hidden confounding and the role of sensitivity analysis as background for Chapter~\ref{chap:confounding-test}.}

\TODO{Motivate why randomized evidence can be used to assess whether observational sensitivity bounds are compatible with a proposed confounding strength.}
